%! Algebra 2 — Unit 2: Quadratic Functions — Study Guide (reference sheet, no problems, no key)
%
% A unit-level bookend component, built by shared/unit.mk the same way unit_cover is (no
% Makefile of its own; `make -C unit02 study_guide`). It is a REFERENCE sheet: vocabulary,
% forms, the graphs to know, and each lesson's target misconception stated as a caution.
% There are no problems here and therefore no answer key — the practice test is the problem
% set students work.
%
% No \namedateperiod: the namestrip rule puts a name row on lesson covers and unit tests only,
% and this is neither — it is a keep-it-in-the-binder handout.
%
% TWO PAGES at 10pt is a hard ceiling. If an edit spills to page 3, trim wording in the tables
% before touching the graph row.
\documentclass[10pt]{article}
\usepackage{algebra2-article}
\usepackage{algebra2-boxes}

% ── Local layout macros ───────────────────────────────────────────────────────
% A compact lesson divider, the same device the practice test uses for its parts
% (headlinebox takes one argument: the background colour).
\newcommand{\sghead}[2]{%
  \boxguard[10]%
  \vspace{4pt}\par\noindent
  \colorbox{forest}{\parbox{\dimexpr\linewidth-2\fboxsep}{%
    \small\color{white}\bfseries #1 \hfill \textmd{\itshape\scriptsize\color{forestmid} #2}}}%
  \vspace{2pt}\par}

% The target misconception for a lesson, stated as a caution.
\newcommand{\sgwatch}[1]{%
  \vspace{1pt}\par\noindent
  \colorbox{redbg}{\parbox{\dimexpr\linewidth-2\fboxsep}{%
    \footnotesize\textbf{\textcolor{redacc}{Watch out:}} #1}}%
  \vspace{1pt}\par}

% Two-column reference table: idea/term on the left, the statement on the right.
% NOT wrapped in a \newenvironment — tabularx scans its body for a literal
% \end{tabularx}, and hiding that token inside an environment's end code fails with
% "File ended while scanning use of \TX@get@body". The spec is a macro instead.
% The left column is a \newcolumntype, not a macro: array does not expand an ordinary
% macro inside a preamble ("Illegal pream-token").
\newcolumntype{G}{>{\bfseries\color{forest}}p{0.245\linewidth}}
\newcommand{\sgtabsetup}{\renewcommand{\arraystretch}{1.06}\footnotesize}

\begin{document}

\pageheader{Unit 2: Quadratic Functions}{Study Guide --- Reference Sheet}

\vspace{-0.30in}
\noindent\footnotesize Everything Unit 2 asks you to \emph{know} --- vocabulary, forms, graphs,
and the six places students most often go wrong. No problems: work the \textbf{practice test}
for those, then come back here for anything you could not explain.\par\normalsize

% ======================================================================
\sghead{2.0 \quad Characteristics of Quadratic Functions}{A2.F.2a, c, d, f, g}
{\sgtabsetup
\begin{tabularx}{\linewidth}{@{} G X @{}}
Parabola & The graph of $f(x)=ax^2+bx+c$, $a\ne0$. It opens \emph{up} when $a>0$ and
  \emph{down} when $a<0$. \\
Vertex & The one turning point, written as a point $(h,k)$: the lowest point if it opens up (a
  \textbf{minimum} point), the highest if it opens down (a \textbf{maximum} point). \\
Max / min value & The \emph{output} at the vertex, $k$ --- a $y$-value, with units. The
  $x$-coordinate $h$ only says \emph{where} it occurs. \\
Axis of symmetry & The line $x=h$ through the vertex. Inputs equally far either side have equal
  outputs, so averaging the two zeros finds the axis. \\
Zeros and $y$-intercept & Zeros ($x$-intercepts) are the inputs where $f(x)=0$: two, one, or
  none --- never three, because the curve turns once. The $y$-intercept is $(0,f(0))$. \\
Increasing / decreasing & Split \emph{at} the axis, given as intervals of $x$. Opening down:
  increasing on $(-\infty,h]$, decreasing on $[h,\infty)$. \\
Domain, range, ends & Domain all reals; range $y\ge k$ (opens up) or $y\le k$ (opens down).
  $a>0$ both ends rise, $a<0$ both ends fall --- always the \emph{same} way. \\
\end{tabularx}}
\sgwatch{The maximum or minimum \emph{value} is an \emph{output}; \emph{where} it happens is an
\emph{input}. A ball's vertex $(1,64)$ means the greatest height is $64$~ft, reached at $t=1$~s
--- ``the maximum is $1$'' answers the wrong question.}

% ======================================================================
\sghead{2.1 \quad Graphing Quadratic Functions \& Transformations}{A.F.2b--d; A2.F.1; A2.F.2a, d}
{\sgtabsetup
\begin{tabularx}{\linewidth}{@{} G X @{}}
Parent function & $y=x^2$, vertex $(0,0)$ --- every quadratic is the parent transformed. \\
Vertex form & $y=a(x-h)^2+k$: vertex $(h,k)$, axis $x=h$; right $h$, up $k$. $(x+2)^2$ is
  $(x-(-2))^2$, so it moves \emph{left} $2$. \\
Sign of $a$ & $a<0$ reflects: opens down, the vertex is a \textbf{maximum}, range $y\le k$.
  $a>0$: a \textbf{minimum}, range $y\ge k$. \\
Size of $a$ & $|a|>1$ narrower (vertical stretch); $0<|a|<1$ wider (vertical compression). Read
  $a$ first (shape, direction), then $h$ and $k$ (position). \\
Standard form & $y=ax^2+bx+c$: $y$-intercept $(0,c)$; axis $x=-\frac{b}{2a}$; substitute that
  $x$ to get the vertex. \\
Intercept form & $y=a(x-p)(x-q)$: zeros $p$ and $q$; the axis is halfway, $x=\frac{p+q}{2}$. \\
One curve, three forms & $(x-1)^2-4=x^2-2x-3=(x+1)(x-3)$: vertex $(1,-4)$, $y$-intercept $-3$,
  zeros $-1$ and $3$. Each form hands over one feature for free. \\
\end{tabularx}}
\sgwatch{The sign in the bracket flips. $(x-1)^2$ moves the parent \emph{right} $1$, and
$(x+1)(x-3)$ has zeros $-1$ and $3$ --- set each bracket equal to zero; never copy the printed
sign.}

% ======================================================================
\sghead{2.2 \quad Solving Quadratics by Factoring}{A2.EO.3b, d; A2.EI.2b, d}
{\sgtabsetup
\begin{tabularx}{\linewidth}{@{} G X @{}}
Factor / GCF & To factor is to write a polynomial as a product. Take out the \textbf{greatest
  common factor} first: $2x^2-8x=2x(x-4)$. \\
Product--sum & For $x^2+bx+c$, find integers $p,q$ with $pq=c$ and $p+q=b$; then
  $x^2+bx+c=(x+p)(x+q)$. \\
Difference of squares & $x^2-k^2=(x+k)(x-k)$ --- the pair is $k,-k$ ($b=0$). A \emph{sum} of
  squares $x^2+k^2$ does not factor. \\
Perfect-square trinomial & $x^2+2kx+k^2=(x+k)^2$ --- a \textbf{double root}: the graph touches
  the axis once. \\
Grouping & For $ax^2+bx+c$, $a\ne1$: two integers with product $a\cdot c$ and sum $b$ split $bx$;
  group in pairs, take out each GCF, then the common binomial. \\
Zero Product Property & If $A\cdot B=0$, then $A=0$ or $B=0$. So: get it $=0$, factor, set each
  factor $=0$. \\
Roots & The solutions of $ax^2+bx+c=0$ are the $x$-intercepts of $y=ax^2+bx+c$. \\
\end{tabularx}}
\sgwatch{The Zero Product Property needs \emph{zero}. Never divide both sides of $x^2=5x$ by $x$
--- that loses $x=0$. Write $x^2-5x=0$, $x(x-5)=0$, so $x=0$ or $5$. And $(x-1)(x+2)=4$ does
\emph{not} mean each factor equals $4$.}

% ======================================================================
\sghead{2.3 \quad Solving by Square Roots \& Complex Numbers}{A2.EI.2b, d; A2.EO.4a--c}
{\sgtabsetup
\begin{tabularx}{\linewidth}{@{} G X @{}}
Square Root Property & If $x^2=k$ and $k\ge0$, then $x=\pm\sqrt{k}$; $(x-h)^2=k$ gives
  $x=h\pm\sqrt{k}$. Isolate the square \emph{first}, and keep the $\pm$ --- two roots. \\
Simplest radical form & No perfect-square factor left under the root: $\sqrt{12}=2\sqrt3$. \\
Imaginary unit & $i=\sqrt{-1}$, so $i^2=-1$. For $k>0$, $\sqrt{-k}=i\sqrt{k}$. $x^2=-4$ has
  roots $\pm2i$ --- which is why $y=x^2+4$ misses the axis. \\
Complex number & $a+bi$ (standard form): real part $a$, imaginary part $b$. $b=0$: real;
  $a=0$, $b\ne0$: pure imaginary. \\
Add / subtract & Combine the real parts and the imaginary parts: $(3+2i)+(1-5i)=4-3i$. \\
Multiply & Distribute, then replace every $i^2$ with $-1$: $(2+i)(3-i)=6+i-i^2=7+i$. \\
Conjugates & $a+bi$ and $a-bi$; their product $a^2+b^2$ is always real: $(2+3i)(2-3i)=13$. \\
\end{tabularx}}
\sgwatch{$i^2=-1$, so pull the $i$ out of each negative radical \emph{before} you multiply:
$\sqrt{-4}\cdot\sqrt{-9}=2i\cdot3i=6i^2=-6$, not $\sqrt{36}=6$. And never leave an $i^2$ in a
final answer.}

% ======================================================================
\sghead{2.4 \quad Completing the Square, the Quadratic Formula \& the Discriminant}{A2.EI.2a, b, d; A2.F.2d}
{\sgtabsetup
\begin{tabularx}{\linewidth}{@{} G X @{}}
Completing the square & For $x^2+bx+c=0$: move $c$, add $\left(\frac b2\right)^2$ to
  \textbf{both} sides, rewrite as $\left(x+\frac b2\right)^2$, square-root with $\pm$. \\
Perfect square, again & $x^2+bx+\left(\frac b2\right)^2=\left(x+\frac b2\right)^2$ --- 2.2's
  pattern run backwards, and the road to 2.1's vertex form. \\
Standard form & $ax^2+bx+c=0$ --- everything on one side \emph{first}; read $a$, $b$, $c$ with
  their signs. \\
Quadratic formula & $x=\dfrac{-b\pm\sqrt{b^2-4ac}}{2a}$ --- completing the square done once, on
  the letters. It solves \emph{every} quadratic; simplify to radical or $p\pm qi$ form. \\
Discriminant & $D=b^2-4ac$, the number under the radical. Square all of $b$: $(-3)^2=9$, not
  $-3^2$. \\
Reading $D$ & $D>0$: two real solutions, two $x$-intercepts. $D=0$: one repeated real, the graph
  touches once. $D<0$: two complex solutions $p\pm qi$, no $x$-intercepts. \\
\end{tabularx}}
\sgwatch{A negative discriminant does \emph{not} mean ``no solution.'' It means no \emph{real}
solutions and no $x$-intercepts --- there are still \textbf{two complex} solutions: $x^2+2x+5=0$
has $D=-16$ and $x=-1\pm2i$.}

% ======================================================================
\sghead{2.5 \quad Quadratic Systems \& Modeling}{A2.EI.3a--d; A2.EI.2a, d; A2.F.2d}
{\sgtabsetup
\begin{tabularx}{\linewidth}{@{} G X @{}}
Solution of a system & An ordered pair $(x,y)$ that makes \emph{both} equations true --- a point
  where the graphs intersect. It needs both coordinates. \\
Substitution & Set the $y$'s equal, move everything to one side, solve the quadratic (any 2.2--2.4
  method), then put each $x$ into the \emph{line} for its $y$. Check in both. \\
Number of solutions & A line and a parabola \textbf{cross} ($2$), \textbf{touch} ($1$), or
  \textbf{miss} ($0$) --- the collapsed equation's $D>0$, $D=0$, $D<0$. \\
Feature $\to$ question & $y$-intercept $h(0)$: the start. Positive zero: when it lands (throw out
  $t<0$). Vertex input $-\frac{b}{2a}$: \emph{when} it peaks. Vertex output: \emph{how high}. \\
Max / min value & The vertex's \textbf{output}: find $-\frac{b}{2a}$, then evaluate there. Answer
  with units. \\
\end{tabularx}}
\sgwatch{The maximum \emph{value} is the vertex's \emph{output}; \emph{when} it happens is the
input. For $h(t)=-16t^2+32t+48$, ``how high?'' is $64$~ft --- not $1$ (the time), not $48$ (the
start), not $3$ (the landing).}

% ======================================================================
\sghead{The six graphs to know on sight}{read each one before the test}
\vspace{2pt}
\noindent
\begin{minipage}[t]{0.16\linewidth}\centering
\begin{tikzpicture}[scale=0.42]
  \draw[step=1, linegray!40, very thin] (-1.6,-0.6) grid (3.6,4.6);
  \draw[->, semithick, charcoal] (-1.6,0)--(3.6,0);
  \draw[->, semithick, charcoal] (0,-0.6)--(0,4.6);
  \begin{scope}
    \clip (-1.6,-0.6) rectangle (3.6,4.6);
    \draw[thick, dashed, slate, domain=-1.6:1.6, samples=40, smooth]
      plot (\x, {(\x)*(\x)});
    \draw[very thick, forest, domain=-0.4:3.6, samples=50, smooth]
      plot (\x, {((\x)-2)*((\x)-2)+1});
  \end{scope}
  \fill[forest] (2,1) circle (3pt);
\end{tikzpicture}\\[1pt]
{\scriptsize\textbf{Shift.} $y=(x-2)^2+1$: the parent (dashed) moved right $2$, up $1$.}
\end{minipage}%
\hfill
\begin{minipage}[t]{0.16\linewidth}\centering
\begin{tikzpicture}[scale=0.42]
  \draw[step=1, linegray!40, very thin] (-1.6,-1.6) grid (3.6,3.6);
  \draw[->, semithick, charcoal] (-1.6,0)--(3.6,0);
  \draw[->, semithick, charcoal] (0,-1.6)--(0,3.6);
  \draw[dashed, slate] (1,-1.6)--(1,3.6);
  \begin{scope}
    \clip (-1.6,-1.6) rectangle (3.6,3.6);
    \draw[very thick, forest, domain=-1.6:3.6, samples=50, smooth]
      plot (\x, {3-((\x)-1)*((\x)-1)});
  \end{scope}
  \fill[forest] (1,3) circle (3pt);
\end{tikzpicture}\\[1pt]
{\scriptsize\textbf{Opens down.} $a<0$: vertex $(1,3)$ is a maximum; range $y\le3$.}
\end{minipage}%
\hfill
\begin{minipage}[t]{0.16\linewidth}\centering
\begin{tikzpicture}[scale=0.42]
  \draw[step=1, linegray!40, very thin] (-1.6,-4.4) grid (3.6,0.8);
  \draw[->, semithick, charcoal] (-1.6,0)--(3.6,0);
  \draw[->, semithick, charcoal] (0,-4.4)--(0,0.8);
  \begin{scope}
    \clip (-1.6,-4.4) rectangle (3.6,0.8);
    \draw[very thick, forest, domain=-1.6:3.6, samples=50, smooth]
      plot (\x, {(\x)*(\x)-2*(\x)-3});
  \end{scope}
  \fill[forest] (-1,0) circle (3pt);
  \fill[forest] (3,0) circle (3pt);
\end{tikzpicture}\\[1pt]
{\scriptsize\textbf{$D>0$.} $y=x^2-2x-3$: two real zeros, $-1$ and $3$.}
\end{minipage}%
\hfill
\begin{minipage}[t]{0.16\linewidth}\centering
\begin{tikzpicture}[scale=0.42]
  \draw[step=1, linegray!40, very thin] (-1.6,-0.6) grid (3.6,4.6);
  \draw[->, semithick, charcoal] (-1.6,0)--(3.6,0);
  \draw[->, semithick, charcoal] (0,-0.6)--(0,4.6);
  \begin{scope}
    \clip (-1.6,-0.6) rectangle (3.6,4.6);
    \draw[very thick, forest, domain=-1.6:3.6, samples=50, smooth]
      plot (\x, {((\x)-1)*((\x)-1)});
  \end{scope}
  \fill[forest] (1,0) circle (3pt);
\end{tikzpicture}\\[1pt]
{\scriptsize\textbf{$D=0$.} $y=(x-1)^2$: one repeated root; it touches the axis at $1$.}
\end{minipage}%
\hfill
\begin{minipage}[t]{0.16\linewidth}\centering
\begin{tikzpicture}[scale=0.42]
  \draw[step=1, linegray!40, very thin] (-1.6,-0.6) grid (3.6,4.6);
  \draw[->, semithick, charcoal] (-1.6,0)--(3.6,0);
  \draw[->, semithick, charcoal] (0,-0.6)--(0,4.6);
  \begin{scope}
    \clip (-1.6,-0.6) rectangle (3.6,4.6);
    \draw[very thick, forest, domain=-1.6:3.6, samples=50, smooth]
      plot (\x, {((\x)-1)*((\x)-1)+1});
  \end{scope}
  \fill[forest] (1,1) circle (3pt);
\end{tikzpicture}\\[1pt]
{\scriptsize\textbf{$D<0$.} $y=(x-1)^2+1$ misses the axis: no real zeros, two complex, $1\pm i$.}
\end{minipage}%
\hfill
\begin{minipage}[t]{0.16\linewidth}\centering
\begin{tikzpicture}[scale=0.42]
  \draw[step=1, linegray!40, very thin] (-2.6,-0.6) grid (2.6,4.6);
  \draw[->, semithick, charcoal] (-2.6,0)--(2.6,0);
  \draw[->, semithick, charcoal] (0,-0.6)--(0,4.6);
  \begin{scope}
    \clip (-2.6,-0.6) rectangle (2.6,4.6);
    \draw[very thick, forest, domain=-2.6:2.6, samples=50, smooth]
      plot (\x, {(\x)*(\x)});
    \draw[very thick, navy] (-2.6,-0.6)--(2.6,4.6);
  \end{scope}
  \fill[redacc] (-1,1) circle (3pt);
  \fill[redacc] (2,4) circle (3pt);
\end{tikzpicture}\\[1pt]
{\scriptsize\textbf{A system.} $y=x^2$ and $y=x+2$ cross twice: $(-1,1)$ and $(2,4)$.}
\end{minipage}

\end{document}
